\chapter{Conclusions}
\label{chap:conclusions}

This project studied two iterative algorithms for solving the LASSO problem arising
in the training of Extreme Learning Machines:
\begin{itemize}[nosep]
    \item \textbf{Algorithm A1 — IRLS}: an instance of the Majorization-Minimization
    framework that replaces the non-smooth $\ell_1$ penalty with a sequence of smooth
    quadratic surrogates, solved via the weighted normal equations.
    \item \textbf{Algorithm A2 — DSM}: a deflected subgradient method with a
    target-level Polyak stepsize, which directly handles non-smoothness without
    any smoothing.
\end{itemize}

\paragraph{Theoretical findings.}
The MM interpretation of IRLS guarantees strict monotone decrease of the objective
at every iteration and, under mild conditions, convergence to a stationary point
that satisfies the $\ell_1$ optimality conditions.
In practice, convergence is linear (geometric), making IRLS very efficient for
small to moderate hidden layer sizes.

DSM provides convergence guarantees of a different kind: the record value
$\bar{f}^{\,i}$ converges to $f^*$ at a sublinear $O(1/\varepsilon^2)$ rate,
which is information-theoretically optimal for non-smooth convex problems.
The deflection mechanism and target-level stepsize significantly improve over the
plain subgradient method, but cannot overcome the fundamental sublinear barrier.

\paragraph{Experimental findings.}
Numerical experiments on synthetic ELM problems confirm all theoretical predictions:
\begin{itemize}
    \item IRLS converges linearly and reaches high accuracy ($f - f^* < 10^{-8}$)
    in fewer than 50 iterations, while DSM requires thousands of iterations for
    comparable accuracy.
    \item Per-iteration, DSM is cheaper ($O(MH)$ vs.\ $O(H^3)$), but the iteration
    count advantage of IRLS dominates for $H \lesssim 1000$.
    \item Scalability experiments show that for $H \gtrsim 1000$ the Cholesky cost
    of IRLS grows as $O(H^3)$, and DSM's per-iteration advantage starts to be
    competitive.
    \item IRLS produces exact sparsity naturally, whereas DSM requires post-processing
    thresholding.
    \item DSM is sensitive to the choice of $\delta_0$ and $\rho$; IRLS is robust to
    its single additional hyperparameter $\varepsilon_{\mathrm{thr}}$.
\end{itemize}

\paragraph{Practical recommendation.}
For the typical ELM setting ($M \gg H$, offline training, high accuracy required),
\textbf{IRLS is the method of choice}.
DSM serves as a useful baseline and becomes competitive only for very large hidden
layers ($H \gtrsim 1000$--$3000$) where $O(H^3)$ Cholesky factorisation is
prohibitively expensive, or in memory-constrained settings where storing
$\mathbf{X}^T\mathbf{X} \in \mathbb{R}^{H \times H}$ is impractical.

\paragraph{Limitations and future work.}
The experiments are limited to synthetic regression datasets with Gaussian features.
Extending the comparison to real-world datasets (e.g.\ image or time-series regression)
would provide additional evidence.
For very large $H$, replacing the Cholesky inner solver in IRLS with Conjugate
Gradient (already implemented) could extend IRLS's competitive range significantly.
Finally, combining the two algorithms — e.g.\ warm-starting DSM from an IRLS solution,
or using IRLS as a preconditioner for DSM — is a natural direction for further work.
